\section{Resumen}

\subsection{Propósito y alcance del estudio}

Este estudio compara dos estructuras de datos para resolver el problema de
\textit{Range Minimum Query} (RMQ): \textbf{Sparse Table} y
\textbf{Segment Tree}. Dado un arreglo de enteros \(A\), una consulta sobre
el rango inclusivo \(A[L..R]\) busca obtener el \textbf{índice del elemento
mínimo}, no solamente su valor. Se consideran un \textbf{escenario estático},
en el que el arreglo permanece fijo, y un \textbf{escenario dinámico},
en el que sus elementos reciben modificaciones puntuales.

El diseño experimental contempla cuatro comparaciones:
\begin{itemize}
    \item \textbf{Construcción en el escenario estático:} medir el tiempo
    necesario para construir ambas estructuras a partir del mismo arreglo,
    considerando diferentes cantidades de elementos.
    \item \textbf{Consultas en el escenario estático:} medir el tiempo de
    respuesta sobre las estructuras ya construidas, utilizando rangos
    aleatorios de distintas longitudes y arreglos de diferentes tamaños.
    \item \textbf{Mantenimiento en el escenario dinámico:} comparar el
    tiempo de una actualización puntual en Segment Tree con el tiempo de
    reconstruir Sparse Table después del mismo cambio. Esta comparación
    representa el costo de mantener respuestas vigentes, no dos operaciones
    de actualización equivalentes.
    \item \textbf{Consultas en el escenario dinámico:} medir el tiempo de
    respuesta después de aplicar los cambios y mantener ambas estructuras
    actualizadas, separándolo del costo de mantenimiento.
\end{itemize}

En cada comparación se utilizarán las mismas entradas y consultas para ambas
estructuras. Además, se realizará un \textbf{análisis del espacio utilizado},
considerando las asignaciones de memoria de los constructores y las reservas
temporales relevantes durante la construcción o reconstrucción. Los resultados
se interpretarán junto con las complejidades teóricas para discutir las
ventajas y limitaciones de cada alternativa. El estudio también abordará su
adaptación a otras operaciones binarias y las propiedades algebraicas que
permiten conservar sus garantías.

\subsection{Estructuras y algoritmos estudiados}

\begin{itemize}
    \item \textbf{Sparse Table:} preprocesa el arreglo almacenando resultados
    para bloques de longitud potencia de dos. Para RMQ permite responder
    consultas en tiempo constante mediante la combinación de dos bloques
    que pueden solaparse. La implementación estudiada no ofrece
    actualizaciones, por lo que en el escenario dinámico se reconstruirá
    después de cada modificación.
    \item \textbf{Segment Tree:} organiza los mínimos de intervalos en una
    estructura jerárquica. Responde consultas y permite actualizar un
    elemento en tiempo logarítmico, recalculando los nodos afectados sin
    reconstruir toda la estructura.
\end{itemize}

Las implementaciones incorporadas proceden del repositorio oficial del curso,
con autoría original de LELE \cite{edaa-sparse-table,edaa-segment-tree}.
